\chapter{Marco Teórico}

\section{Inteligencia artificial en juegos}

Los juegos han sido utilizados de manera recurrente como dominio de prueba para
la inteligencia artificial porque ofrecen entornos con reglas explícitas,
objetivos medibles y resultados verificables. Desde los primeros trabajos de
Arthur Samuel sobre aprendizaje automático aplicado al juego de damas
\cite{Samuel1959}, los juegos han permitido estudiar problemas de búsqueda,
planeación, toma de decisiones, aprendizaje y evaluación de estrategias.

Una razón por la cual los juegos son relevantes para la inteligencia artificial
es que concentran, en un entorno controlado, problemas que también aparecen en
dominios más amplios: seleccionar acciones bajo restricciones, anticipar
consecuencias, competir con otros agentes, administrar recursos computacionales y
tomar decisiones sin explorar por completo todas las posibilidades. En juegos
simples, como Gato, es posible recorrer una parte importante del espacio de
estados. En juegos más complejos, como Ajedrez, el número de estados y secuencias
posibles crece de forma combinatoria, por lo que se requieren estrategias de
búsqueda y evaluación.

Históricamente, muchos sistemas exitosos de inteligencia artificial para juegos
han sido motores especializados. Estos sistemas alcanzan alto desempeño porque
incorporan conocimiento particular del dominio: reglas programadas para un juego,
funciones de evaluación ajustadas a sus piezas o posiciones, representaciones
hechas a la medida y optimizaciones diseñadas para un tablero específico. Aunque
este enfoque puede producir agentes fuertes en un juego determinado, también
genera arquitecturas rígidas. Un programa diseñado para jugar Ajedrez no puede,
por ese solo hecho, jugar Damas, Go o Gato. La limitación no se encuentra en su
capacidad de cálculo, sino en que su conocimiento está unido al juego para el que
fue construido.

El problema anterior motiva el estudio de agentes capaces de operar sobre más de
un juego sin reprogramar su motor. Esta línea de investigación se conoce como
General Game Playing y busca mover parte del trabajo intelectual desde el
programador hacia el agente: el sistema no debe depender de que el desarrollador
codifique manualmente cada regla, sino de su capacidad para interpretar una
descripción formal del juego y razonar sobre ella.

\section{Juegos de información perfecta}

Un juego de información perfecta es aquel en el que todos los jugadores conocen
el estado completo del juego cuando toman una decisión. No existen cartas ocultas,
información privada, azar no observado ni reglas desconocidas para alguno de los
participantes. Ejemplos clásicos de este tipo de juegos son Gato, Ajedrez y
Damas. En ellos, el jugador puede ver la configuración actual y determinar qué
acciones son legales de acuerdo con las reglas.

Para estudiar estos juegos desde una perspectiva computacional conviene
representarlos mediante componentes formales. De manera general, un juego puede
describirse a partir de un conjunto de jugadores, un conjunto de estados, un
estado inicial, acciones disponibles, una función de transición, estados
terminales y valores de utilidad o recompensa. Esta formalización permite que el
agente razone con conceptos abstractos en lugar de depender de nombres concretos
como ``reina'', ``torre'', ``X'' u ``O''.

En esta representación, un estado contiene toda la información necesaria para
describir una situación del juego. Una acción legal es un movimiento permitido en
ese estado. La transición indica cómo cambia el juego después de aplicar una
acción. Un estado terminal representa el final de la partida y la utilidad indica
qué tan favorable fue el resultado para cada jugador. Esta separación conceptual
es fundamental para un General Game Player, porque permite que el motor de
decisión opere sobre una interfaz común aun cuando cambien las reglas del juego.

\section{General Game Playing}

Un General Game Player (GGP) es un agente que recibe una descripción formal de un
juego y debe jugarlo sin intervención humana ni conocimiento específico
preprogramado para ese dominio. Genesereth, Love y Pell definen el área como el
estudio de sistemas capaces de aceptar la descripción formal de un juego y jugar
efectivamente sin depender de algoritmos diseñados de antemano para juegos
particulares \cite{Genesereth2005GGP}.

La diferencia principal entre un motor especializado y un General Game Player se
encuentra en la ubicación del conocimiento del juego. En un motor especializado,
las reglas y heurísticas suelen estar dentro del código fuente. En un GGP, las
reglas se proporcionan como entrada externa. El agente debe interpretarlas,
construir una representación interna y usar un proceso general de decisión. Por
ello, la generalidad no significa únicamente que el sistema pueda ejecutar dos
juegos, sino que el mismo núcleo algorítmico pueda operar al cambiar la
descripción del juego.

Este enfoque tiene una implicación directa para el presente Trabajo Terminal: el
objetivo no consiste en construir un motor competitivo de Ajedrez comparable con
sistemas especializados, sino en diseñar una arquitectura desacoplada. La
validación debe demostrar que el sistema puede cargar reglas externas, generar
movimientos legales, reconocer estados terminales y seleccionar jugadas bajo un
criterio general de búsqueda. La prueba de generalidad se encuentra en cambiar
el archivo de reglas sin modificar el motor del agente.

\section{Modelo formal de juego}

El modelo formal de juego en GGP puede entenderse como una máquina de estados. De
acuerdo con la descripción general presentada por Genesereth, Love y Pell, un
juego finito puede modelarse mediante estados, roles, acciones legales, función
de actualización, estado inicial, condiciones terminales y valores de objetivo
\cite{Genesereth2005GGP}. Esta visión permite que el juego sea procesado por una
computadora como un sistema dinámico.

El conjunto de estados representa todas las configuraciones posibles. Los roles
corresponden a los jugadores. Para cada estado, las reglas determinan qué
acciones puede ejecutar cada rol. La función de transición toma el estado actual
y las acciones realizadas para producir un nuevo estado. El juego continúa hasta
alcanzar un estado terminal. En ese punto, se asigna una recompensa a cada
jugador.

Para que esta abstracción sea útil en un sistema computacional, el motor del
agente debe poder realizar consultas generales. Por ejemplo, necesita obtener el
estado inicial, listar acciones legales, aplicar una acción, verificar si un
estado es terminal y consultar la recompensa. Estas operaciones pueden expresarse
conceptualmente como:

\begin{itemize}
    \item obtener el estado inicial del juego;
    \item consultar las acciones legales para un jugador en un estado;
    \item aplicar una acción y producir un estado sucesor;
    \item determinar si un estado es terminal;
    \item calcular la utilidad o recompensa para cada jugador.
\end{itemize}

La importancia de esta interfaz es que el motor de búsqueda no necesita conocer
las reglas concretas del dominio. En lugar de preguntar si un caballo puede
moverse en forma de ``L'' o si una casilla de Gato está vacía, pregunta qué
acciones son legales desde el estado actual. El conocimiento específico debe
estar en la descripción externa y en la representación interna derivada de ella,
no en el algoritmo de decisión.

\section{Lenguaje de descripción de juegos}

Un lenguaje de descripción de juegos permite expresar formalmente las reglas de
un juego para que puedan ser interpretadas por un agente. En General Game
Playing, el lenguaje clásico es Game Description Language (GDL). GDL fue
diseñado para describir juegos discretos mediante reglas lógicas, de modo que un
agente pueda conocer el estado inicial, las acciones legales, las transiciones,
las condiciones terminales y las recompensas \cite{Love2008GDL}.

GDL utiliza predicados especiales para representar los componentes del juego.
Entre los más importantes se encuentran:

\begin{itemize}
    \item \texttt{role}, para declarar jugadores o roles;
    \item \texttt{init}, para indicar hechos verdaderos en el estado inicial;
    \item \texttt{true}, para consultar hechos del estado actual;
    \item \texttt{legal}, para definir acciones permitidas;
    \item \texttt{does}, para representar acciones ejecutadas;
    \item \texttt{next}, para describir hechos del estado siguiente;
    \item \texttt{terminal}, para establecer condiciones de fin de juego;
    \item \texttt{goal}, para asignar valores de recompensa.
\end{itemize}

Esta forma de representación muestra que GDL no funciona como un archivo simple
de configuración. No basta con indicar el tamaño del tablero o el número de
jugadores. El lenguaje debe expresar relaciones, condiciones, variables y reglas
de inferencia. Por ejemplo, una acción puede ser legal solo si una casilla está
vacía y si corresponde el turno del jugador. Esta capacidad declarativa hace que
GDL sea potente, pero también introduce complejidad en el diseño del intérprete.

Cropper, Evans y Law explican que las reglas de GGP se describen como programas
lógicos y que especifican elementos como estado inicial, movimientos legales,
actualización del estado y terminación del juego \cite{Cropper2019IGGP}. Aunque
su trabajo se enfoca en aprender reglas a partir de trazas de juego, resulta
relevante porque evidencia la complejidad lógica de representar juegos de manera
general.

\section{GDL, Ludii y lenguajes de dominio específico}

El uso de GDL clásico no es la única alternativa para describir juegos. Existen
otros lenguajes de descripción que buscan equilibrar expresividad, legibilidad y
facilidad de implementación. Ludii, por ejemplo, propone un lenguaje de
descripción de juegos de mayor nivel, con conceptos más cercanos al diseño de
juegos, como tablero, piezas y movimientos \cite{Soemers2024Ludii}.

La comparación entre GDL y Ludii es importante para este Trabajo Terminal porque
plantea una decisión arquitectónica: implementar GDL completo, adoptar un
subconjunto de GDL o diseñar un lenguaje de dominio específico (DSL) inspirado en
estos enfoques. GDL ofrece una base académica fuerte y una semántica lógica
general, pero puede resultar complejo para un prototipo de alcance acotado.
Ludii, por su parte, muestra que un lenguaje puede conservar generalidad y al
mismo tiempo ser más legible para diseñadores y desarrolladores.

Un DSL propio permitiría concentrarse en los elementos necesarios para el
alcance del proyecto: jugadores, estados, acciones, reglas de legalidad,
transiciones, condiciones terminales y recompensas. Esta opción debe justificarse
formalmente para evitar que el lenguaje se convierta en una colección de casos
especiales. El criterio de diseño debe ser que el lenguaje permita describir
juegos dentro del alcance definido: discretos, deterministas, por turnos y de
información perfecta.

\section{Intérprete y representación interna}

El archivo de reglas no es suficiente por sí mismo para jugar. Se requiere un
intérprete que lea la descripción, valide su estructura y construya una
representación interna consultable por el motor del agente. Este proceso suele
incluir análisis léxico, análisis sintáctico, validación semántica y construcción
de un modelo interno.

El análisis léxico separa el texto en unidades reconocibles. El análisis
sintáctico verifica que esas unidades respeten la gramática del lenguaje. La
validación semántica revisa que las reglas sean coherentes con el modelo del
juego, por ejemplo que las acciones referencien jugadores existentes o que una
condición terminal pueda evaluarse. Finalmente, la representación interna
organiza esa información para responder consultas durante la partida.

La representación interna debe funcionar como puente entre el lenguaje textual y
los algoritmos de decisión. Una vez construido el modelo, el motor no debería
procesar directamente cadenas de texto, sino estructuras que permitan evaluar
legalidad, transiciones y recompensas. Esta separación mejora la mantenibilidad:
en principio, el lenguaje textual podría cambiar sin modificar por completo el
motor de búsqueda, siempre que se preserve la interfaz conceptual del juego.

En implementaciones avanzadas de GGP, las reglas también pueden compilarse en
estructuras más eficientes. Goldwaser y Thielscher mencionan que una descripción
GDL puede convertirse en una red proposicional, o PropNet, con nodos lógicos,
nodos de entrada, salida y transición \cite{Goldwaser2020DRLGGP}. Esta idea es
relevante porque interpretar reglas desde cero en cada nodo de búsqueda puede
ser costoso. Aunque implementar una PropNet completa puede quedar fuera de una
primera versión, el concepto ayuda a entender por qué la eficiencia de la
representación interna es tan importante como el algoritmo de búsqueda.

\section{Árboles de juego y explosión combinatoria}

Una vez que el sistema puede generar acciones legales y aplicar transiciones, se
puede construir un árbol de juego. La raíz del árbol representa el estado actual.
Cada rama representa una acción legal y cada nodo hijo representa el estado
resultante. Al repetir este proceso, el agente puede anticipar posibles futuros y
comparar las consecuencias de distintas decisiones.

El problema es que el número de nodos crece rápidamente. Si un estado tiene un
factor de ramificación promedio de \(b\) acciones y se explora hasta una
profundidad \(d\), el número de nodos puede crecer en el orden de \(b^d\). En
juegos pequeños, este crecimiento puede ser manejable. En juegos complejos,
explorar todas las posibilidades resulta impracticable.

La explosión combinatoria obliga a limitar la búsqueda. El agente debe decidir
cuánto explorar, qué ramas priorizar y cómo evaluar estados no terminales. Aquí
aparecen algoritmos como Minimax, poda alfa-beta y Monte Carlo Tree Search, así
como funciones heurísticas que estiman la calidad de un estado sin llegar
necesariamente al final de la partida.

\section{Minimax}

Minimax es un algoritmo clásico para juegos adversariales de dos jugadores. Su
principio básico consiste en suponer que un jugador intenta maximizar su utilidad
mientras el oponente intenta minimizarla. En un árbol de juego, los niveles se
alternan entre decisiones del jugador MAX y decisiones del jugador MIN.

Si el árbol puede explorarse hasta estados terminales, Minimax puede seleccionar
una acción óptima bajo el supuesto de que ambos jugadores actúan racionalmente.
El valor de un estado terminal se obtiene directamente de la utilidad del juego.
El valor de un estado no terminal se calcula propagando hacia arriba los valores
de sus sucesores: MAX elige el mayor valor disponible y MIN elige el menor.

En la práctica, muchos juegos no permiten explorar hasta el final. Por ello, se
establece una profundidad máxima y se evalúan los estados frontera mediante una
función heurística. Esta adaptación permite aplicar Minimax en juegos más
grandes, pero introduce una pregunta central para GGP: cómo evaluar estados sin
usar conocimiento específico del juego.

\section{Poda alfa-beta}

La poda alfa-beta es una optimización de Minimax. No cambia la decisión final
cuando se aplica bajo las mismas condiciones de evaluación, pero evita explorar
ramas que no pueden modificar el resultado. Para ello mantiene dos valores:
alfa, que representa la mejor alternativa encontrada por MAX, y beta, que
representa la mejor alternativa encontrada por MIN.

Cuando el algoritmo detecta que una rama no puede producir una decisión mejor que
otra ya conocida, deja de explorarla. De esta manera, reduce el número de nodos
visitados y permite alcanzar mayor profundidad con el mismo presupuesto
computacional. La efectividad de la poda depende del orden en que se exploran
las acciones; si las mejores acciones se revisan temprano, se pueden descartar
más ramas.

Para este Trabajo Terminal, alfa-beta es relevante porque permite comparar dos
ideas: la decisión producida por Minimax y la eficiencia lograda al podar ramas.
En una evaluación experimental, métricas como nodos explorados, tiempo de
decisión y profundidad alcanzada pueden mostrar si la optimización mejora el
rendimiento sin alterar la legalidad ni la coherencia de las jugadas.

\section{Funciones heurísticas generales}

Una función heurística estima qué tan favorable es un estado cuando no es posible
explorar el árbol hasta el final. En motores especializados, las heurísticas
suelen incorporar conocimiento del dominio. En Ajedrez, por ejemplo, puede
asignarse valor a piezas o evaluar control del centro. Sin embargo, esas reglas
no son generales. Si el agente conoce el valor de una reina, entonces contiene
conocimiento específico de Ajedrez.

El reto en General Game Playing es construir evaluaciones que no dependan de la
semántica particular de un juego. Una heurística general puede considerar
propiedades estructurales, como movilidad, cantidad de acciones legales,
proximidad a estados terminales, cambios de recompensa o resultados obtenidos
por simulación. Ninguna de estas medidas garantiza por sí sola buen desempeño en
todos los juegos, pero pueden servir como aproximaciones dentro de un agente
agnóstico al dominio.

Este problema es uno de los puntos más importantes del proyecto. La búsqueda
limitada necesita una forma de decidir qué estados parecen mejores, pero la
generalidad exige que esa evaluación no esté escrita para un dominio concreto.
Por ello, el análisis de heurísticas generales y métodos basados en simulación
forma parte central de la fundamentación teórica del motor de decisión.

\section{Monte Carlo Tree Search}

Monte Carlo Tree Search (MCTS) es una familia de métodos de búsqueda que estima
el valor de las acciones mediante simulaciones. En lugar de depender
exclusivamente de una función heurística explícita, MCTS construye
progresivamente un árbol de búsqueda y utiliza resultados de partidas simuladas
para orientar la exploración \cite{Browne2012MCTS}.

El ciclo general de MCTS se compone de cuatro etapas:

\begin{itemize}
    \item selección, donde se baja por el árbol eligiendo nodos prometedores;
    \item expansión, donde se agrega un nuevo nodo;
    \item simulación, donde se juega desde ese nodo hasta obtener un resultado;
    \item retropropagación, donde el resultado actualiza los valores de los
    nodos visitados.
\end{itemize}

MCTS resulta atractivo para GGP porque puede reducir la dependencia de
heurísticas específicas. Si el sistema cuenta con un modelo de juego capaz de
generar acciones legales, aplicar transiciones y detectar terminales, entonces
puede simular partidas y estimar el valor de acciones a partir de resultados.
No obstante, su rendimiento depende de la calidad del modelo interno y del
presupuesto de simulación disponible.

\section{UCT}

Una cuestión central en MCTS es decidir qué rama explorar. Si el agente explora
solo la acción que hasta ahora parece mejor, puede ignorar alternativas poco
visitadas que resultarían superiores. Si explora demasiado al azar, puede perder
tiempo en acciones de baja calidad. UCT, o Upper Confidence Bounds applied to
Trees, aborda este equilibrio entre explotación y exploración.

UCT combina el valor promedio observado de una acción con un término que favorece
acciones menos exploradas. De forma conceptual, una acción es atractiva si ya ha
producido buenos resultados o si aún no ha sido visitada lo suficiente. Este
criterio permite que MCTS refine progresivamente sus estimaciones sin abandonar
por completo la exploración.

En el contexto de General Game Playing, UCT ha sido importante porque ofrece una
estrategia relativamente general para seleccionar acciones en árboles de juego
sin requerir una heurística escrita manualmente para cada dominio. Goldwaser y
Thielscher señalan que agentes exitosos de GGP han usado enfoques basados en
búsqueda y heurísticas generales, y que posteriormente se han explorado métodos
basados en aprendizaje por refuerzo profundo \cite{Goldwaser2020DRLGGP}.

\section{Aprendizaje por refuerzo y aprendizaje profundo}

El aprendizaje por refuerzo estudia agentes que aprenden mediante interacción
con un entorno. En cada paso, el agente observa un estado, ejecuta una acción,
recibe una recompensa y transita a un nuevo estado. A diferencia del aprendizaje
supervisado, no siempre existe una respuesta correcta inmediata para cada
estado; el agente debe aprender a partir de recompensas que pueden aparecer al
final de una secuencia de acciones.

En juegos, el aprendizaje por refuerzo puede utilizar resultados como victoria,
derrota o empate para ajustar una política de decisión. Cuando se combina con
redes neuronales profundas, se habla de aprendizaje por refuerzo profundo. Este
enfoque permite aproximar funciones complejas de valor o política, pero suele
requerir muchos datos, entrenamiento intensivo y representaciones adecuadas del
estado.

Para un GGP, el aprendizaje por refuerzo profundo plantea una dificultad
adicional: el agente no solo debe aprender a jugar, sino operar sobre juegos
descritos externamente. Esto implica que la representación del estado y de las
acciones debe ser suficientemente general. Por esta razón, el aprendizaje
profundo se considera en este trabajo principalmente como antecedente y línea de
investigación avanzada, no como requisito central de implementación.

\section{AlphaZero como antecedente moderno}

AlphaZero representa un antecedente importante porque combinó autojuego,
aprendizaje por refuerzo profundo y búsqueda Monte Carlo para alcanzar alto
desempeño en Ajedrez, Shogi y Go \cite{Silver2018AlphaZero}. Su arquitectura
utiliza una red neuronal para estimar una política de acciones y un valor de
estado, mientras MCTS refina la selección de movimientos mediante búsqueda.

Sin embargo, AlphaZero no resuelve por sí mismo el problema central de este
Trabajo Terminal. Aunque emplea un algoritmo general en más de un juego, requiere
simuladores y representaciones adecuadas para cada dominio. En un GGP, el reto
principal es distinto: interpretar reglas externas, construir el modelo del
juego y permitir que el motor de decisión opere sin conocimiento específico
codificado para cada caso.

Goldwaser y Thielscher extienden ideas de aprendizaje por refuerzo profundo al
contexto de GGP y muestran que existe una línea de investigación que busca unir
descripciones formales de juegos con métodos inspirados en AlphaZero
\cite{Goldwaser2020DRLGGP}. Para el alcance de este Trabajo Terminal, estas
ideas sirven como referencia del estado del arte y ayudan a delimitar el
proyecto: se busca una arquitectura general e interpretable, no replicar un
sistema de entrenamiento profundo a gran escala.

\section{Síntesis del marco teórico}

El desarrollo de un General Game Player requiere integrar varias áreas de la
inteligencia artificial y de las ciencias computacionales. Primero, se necesita
una representación formal del juego que separe estados, acciones, transiciones,
terminalidad y recompensas. Segundo, se requiere un lenguaje de descripción o DSL
capaz de expresar reglas sin codificarlas directamente en el motor. Tercero, es
necesario un intérprete que transforme esa descripción en una representación
interna consultable. Finalmente, el agente debe seleccionar acciones mediante
algoritmos generales de búsqueda o simulación.

La cadena conceptual del proyecto puede resumirse como:

\begin{center}
    reglas externas $\rightarrow$ intérprete $\rightarrow$ representación interna
    $\rightarrow$ acciones legales $\rightarrow$ árbol de juego $\rightarrow$
    decisión.
\end{center}

La aportación del trabajo no radica en construir el mejor jugador de un juego
particular, sino en demostrar una arquitectura desacoplada. El motor debe
mantenerse estable cuando cambia la descripción del juego. Por ello, Gato sirve
como caso funcional completo para validar el flujo de reglas, estados y
terminalidad, mientras Ajedrez funciona como caso controlado de mayor
complejidad. Esta delimitación permite estudiar generalidad operacional sin
confundir el prototipo con un motor especializado de competencia.
